\ExerciseSection

\begin{enumerate}
\def\labelenumi{\roman{enumi})}
\tightlist
\item
  \begin{enumerate}
  \def\labelenumii{\alph{enumii})}
  \tightlist
  \item
    Construct a topological space consisting of four points which satisfies axiom \((\mathbf{O}_{\mathbf{V}})\) but not axiom \((\mathbf{O}_{\mathbf{III}})\) .
  \end{enumerate}
\end{enumerate}

\begin{enumerate}
\def\labelenumi{\alph{enumi})}
\setcounter{enumi}{1}
\item
  A topological space X which satisfies axiom \((\mathrm{O}_{\mathrm{V}})\) also satisfies axiom \((\mathrm{O}_{\mathrm{III}})\) if and only if it has the following property: every closed subset of X is the intersection of its neighbourhoods. X is then uniformizable and the completely regular space associated with X (§ 1, Exercise 4) is normal.
\item
  If a topological space satisfies axioms (C) and \((\mathrm{O}_{\mathrm{III}})\) , it satisfies \((\mathrm{O}_{\mathrm{V}})\) , and the associated completely regular space is compact.
\end{enumerate}

\begin{enumerate}
\def\labelenumi{\arabic{enumi})}
\setcounter{enumi}{1}
\tightlist
\item
  Let X be a topological space which satisfies axiom \((\mathbf{O}_{\mathbf{v}})\) .
\end{enumerate}

\begin{enumerate}
\def\labelenumi{\alph{enumi})}
\item
  Show that the relation \(R: \{x\} \cap \{y\} \neq \emptyset\) between two points x, y of X is equivalent to the relation \(R' : "f(x) = f(y)\) for every real-valued continuous function f on X''; hence show that R is an equivalence relation on X.
\item
  Show that the quotient space \(\mathbf{Y} = \mathbf{X} / \mathbf{R}\) is normal and that, if \(\varphi: \mathbf{X} \to \mathbf{Y}\) is the canonical mapping, every continuous real-valued function \(f\) on \(\mathbf{X}\) can be written in the form \(f = g \circ \varphi\) , where \(g\) is a continuous real-valued function on \(\mathbf{Y}\) .
\end{enumerate}

\begin{enumerate}
\def\labelenumi{\arabic{enumi})}
\setcounter{enumi}{2}
\tightlist
\item
  If X is a topological space, the following conditions are equivalent:
\end{enumerate}

\begin{enumerate}
\def\labelenumi{\alph{enumi})}
\item
  Every subspace of X satisfies axiom \((\mathrm{O}_{\mathrm{V}}^{\prime})\) .
\item
  Every open subspace of X satisfies axiom \((\mathbf{O}_{\mathbf{V}}^{\prime})\) .
\item
  Given any two subsets A and B of X such that \(A \cap \overline{B} = B \cap \overline{A} = \varnothing\) , there exist two disjoint open sets U, V such that \(A \subset U\) and \(B \subset V\) .
\end{enumerate}

A topological space \(\mathbf{X}^{-}\) is said to be completely normal if it is Hausdorff and satisfies these conditions. Every metrizable space is completely normal. A compact space need not be completely normal.

\begin{enumerate}
\def\labelenumi{\arabic{enumi})}
\setcounter{enumi}{3}
\tightlist
\item
  Every linearly ordered set X endowed with either of the topologies \(\mathcal{C}_{+}(\mathbf{X})\) , \(\mathcal{C}_{-}(\mathbf{X})\) (Chapter I, § 2, Exercise 5) is completely normal. (If A and B are two subsets of X such that \(\mathbf{A} \cap \overline{\mathbf{B}} = \mathbf{B} \cap \overline{\mathbf{A}} = \emptyset\) , define a neighbourhood \(V_{x}\) of \(x\) for each \(x \in \mathbf{A}\) , and a neighbourhood \(W_{y}\) of \(y\) for each \(y \in \mathbf{B}\) , such that \(V_{x} \cap W_{y} = \emptyset\) for all \(x \in \mathbf{A}\) and all \(y \in \mathbf{B}\) .)
\end{enumerate}

¶ 5) Every linearly ordered set X, endowed with the topology \(\mathcal{G}_{0}(\mathbf{X})\) (Chapter I, § 2, Exercise 5), is completely normal. {[}Consider first the case where X is compact; show with the help of Chapter IV, § 2, Exercise 6 that every open set in X is the union of mutually disjoint open intervals; use this result to prove the proposition, by considering (in the notation of Exercise 3) the complement of the closed set \(\overline{\mathbf{A}}\cap\overline{\mathbf{B}}\) , and then the complement of \(\overline{\mathbf{B}}\) ; to pass to the general case, in which X is arbitrary, use Chapter IV, § 4, Exercise 7.{]}

\begin{enumerate}
\def\labelenumi{\arabic{enumi})}
\setcounter{enumi}{5}
\item
  Show that, in a normal space X, every subspace Y of X which is a countable union of closed sets is normal. {[}Let A, B be two disjoint closed subsets of Y, and suppose that Y is the union of an ascending sequence \((\mathbf{Y}_n)\) of closed subsets of \(\mathbf{X}\) . Define by induction two sequences \((\mathbf{U}_n), (\mathbf{V}_n)\) of open subsets of \(\mathbf{X}\) such that: (i) \(\overline{\mathbf{U}_n \cap \mathbf{Y}_n}\) and \(\overline{\mathbf{V}_n \cap \mathbf{Y}_n}\) do not intersect; (ii) \(\mathbf{U}_n \cap \mathbf{Y}_n\) contains \(\mathbf{A} \cap \mathbf{Y}_n\) and all the sets \(\overline{\mathbf{U}_i \cap \mathbf{Y}_i}\) for \(1 \leqslant i \leqslant n\) ; and \(\mathbf{V}_n \cap \mathbf{Y}_n\) contains \(\mathbf{B} \cap \mathbf{Y}_n\) and all \(\overline{\mathbf{V}_i \cap \mathbf{Y}_i}, 1 \leqslant i \leqslant n\) .{]}
\item
  A topological space X is said to be perfectly normal if it is normal and if each closed subset of X is a countable intersection of open sets (or, equivalently, if each open subset of X is a countable union of closed sets).
\end{enumerate}

\begin{enumerate}
\def\labelenumi{\alph{enumi})}
\item
  A Hausdorff space X is perfectly normal if and only if, given any closed subset A of X, there exists a real-valued continuous function f on X such that \(\overline{f}^{1}(0)=A\) {[}use the method of Exercise 16 a) of §1{]}. b) Show that every perfectly normal space X is completely normal and that every subspace of X is perfectly normal {[}use Exercises 3 b) and 6{]}.
\item
  Every lower semi-continuous real-valued function on a perfectly normal space is the upper envelope of a sequence of continuous functions (use the method of § 2, no. 7, Proposition 11).
\item
  The compact space X obtained by adjoining a point at infinity to an uncountable discrete space is completely normal but not perfectly normal, and every subspace of X is paracompact.
\end{enumerate}

§ 8) Show that the non-metrizable compact space X defined in Exercise 13 a) of § 2 is perfectly normal, and that the non-metrizable compact space Y defined in Exercise 13 d) of § 2 is completely normal but not perfectly normal.

¶ 9) a) Let X be a paracompact space in which every point has a countable fundamental system of neighbourhoods, and let Y be a countably compact normal space (§ 2, Exercise 14). Show that X × Y is normal. {[}Let A, B be two disjoint closed subsets of X × Y; for each x ∈ X, show that there is a neighbourhood U\_x of x in X and two open sets V\_x, W\_x in Y such that \(\overline{V}_x \cap \overline{W}_x = \emptyset\) and such that for each z ∈ U\_x we have A(z) ⊂ V\_x and B(z) ⊂ W\_x; to prove this, argue by contradiction. Let (T\_λ)\_\{λ ∈ L\} be a locally finite open covering of X, finer than the covering (U\_x)\_\{x ∈ X\}, and let (f\_λ) be a partition of unity subordinate to (T\_λ); for each λ ∈ L, let x\_λ be such that T\_λ ⊂ U\_\{x\_λ\} and let g\_λ be a continuous mapping of Y into {[}0, 1{]} which is equal to 0 on \(\overline{V}_{x_λ}\) and to 1 on \(\overline{W}_{x_λ}\) ; consider the function \(\sum_{\lambda} f_\lambda(x) g_\lambda(y)\) on X × Y.{]}

\begin{enumerate}
\def\labelenumi{\alph{enumi})}
\setcounter{enumi}{1}
\tightlist
\item
  Let Y = {[}a, b{[} be the locally compact space defined in Chapter I, § 9, Exercise 12, and let X = {[}a, b{]} be the compact space obtained by adjoining a point at infinity to Y. Y is normal (Exercise 4) and countably compact (§ 2, Exercise 15), but the product X×Y is not normal (use Chapter II, § 4, Exercise 4).
\end{enumerate}

¶ 10) Let L be an uncountable set and let X denote the completely regular space \(\mathbf{N}^{\mathrm{L}}\) (N carrying the discrete topology). Let A (resp. B) be the set of all \(x = (x(\lambda))_{\lambda \in \mathrm{L}}\) in X such that, for each integer \(k \neq 0\) (resp. \(k \neq 1\) ) the set of \(\lambda \in \mathbf{L}\) such that \(x(\lambda) = k\) has at most one element. a) Show that A and B are disjoint closed subsets of X.

\begin{enumerate}
\def\labelenumi{\alph{enumi})}
\setcounter{enumi}{1}
\item
  Let U, V be two open subsets of X such that \(A \subset U\) and \(B \subset V\) . Let \(\Phi\) denote the set of all elementary sets in X (Chapter I, § 4, no. 1) whose projections which are distinct from N consist of a single element; for each \(W \in \Phi\) , let \(\mathrm{H}(\mathrm{W})\) be the set of all \(\lambda \in L\) such that \(\mathrm{pr}_{\lambda}(\mathrm{W})\) consists of a single element. Show that there is a sequence \((x_{n})\) of points of A, a sequence \((\mathbf{U}_{n})\) of sets of \(\Phi\) , a sequence \((\lambda_{n})\) of distinct elements of L and a strictly increasing sequence \((m(n))_{n \in \mathbb{N}}\) of integers with the following properties: (i) \(U_{n} \subset U\) is a neighbourhood of \(x_{n}\) ; (ii) \(\mathrm{H}(\mathrm{U}_{n})\) is the set of all \(\lambda_{k}\) such that \(k \leqslant m(n)\) ; (iii) \(x_{0}(\lambda) = 0\) for each \(\lambda \in L\) , and for each n \textgreater{} 0, \(x_{n}(\lambda_{k}) = k\) if \(k \leqslant m(n - 1)\) and \(x_{n}(\lambda) = 0\) for all other \(\lambda \in L\) .
\item
  Let \(y \in \mathbf{B}\) be the point such that \(y(\lambda_k) = k\) for each integer \(k\) and \(y(\lambda) = \mathbf{1}\) for all \(\lambda \in \mathbf{L}\) other than the \(\lambda_k\) . Let \(\mathbf{V}_0 \subset \mathbf{V}\) be a set of \(\Phi\) which contains \(y\) , and let \(n\) be an integer such that \(\lambda_k \in \mathbf{L} - \mathbf{H}(\mathbf{V}_0)\) for all \(k > m(n)\) . Show that \(\mathbf{U}_{n+1} \cap \mathbf{V}_0 \neq \emptyset\) , and deduce that \(\mathbf{X}\) is not normal.
\end{enumerate}

\begin{enumerate}
\def\labelenumi{\arabic{enumi})}
\setcounter{enumi}{10}
\tightlist
\item
  Deduce from Exercise 10 that:
\end{enumerate}

\begin{enumerate}
\def\labelenumi{\alph{enumi})}
\item
  If a product \(\prod_{i\in I}X_i\) of non-empty Hausdorff spaces is normal, then \(X_i\) is countably compact (§ 2, Exercise 14) for all save a countable set of indices.
\item
  A product \(\prod_{i\in I}X_i\) of non-empty metrizable spaces is normal if and only if \(X_{i}\) is compact for all but a countable set of indices; and the product space \(\prod_{i\in I}X_{i}\) is then paracompact.
\end{enumerate}

\begin{enumerate}
\def\labelenumi{\arabic{enumi})}
\setcounter{enumi}{11}
\tightlist
\item
  Let X, Y be two Hausdorff spaces.
\end{enumerate}

\begin{enumerate}
\def\labelenumi{\alph{enumi})}
\item
  Suppose that there is a closed set A in X which is not a countable intersection of open sets, and that there is a countably infinite subset B of Y which is not closed. Let \(b \in \overline{B} - B\) , and consider the subsets \(C = A \times B\) , \(D = (X - A) \times \{b\}\) in \(X \times Y\) . Show that \(C \cap \overline{D} = \overline{C} \cap D = \varnothing\) , but that there is no pair U, V of open subsets of \(X \times Y\) such that \(C \subset U\) , \(D \subset V\) and \(U \cap V = \varnothing\) .
\item
  If \(X \times Y\) is completely normal then either (i) one of the spaces X, Y is perfectly normal, or (ii) one of the spaces X, Y has the property that every countably infinite subset is closed (cf.~§ 5, Exercise 16).
\item
  Let \(X = [a, b]\) be an uncountable well-ordered set such that \([a, x]\) is countable for each x \textless{} b. Endow X with the topology in which \(\{x\}\) is open for each x \textless{} b and the intervals \([x, b]\) (x \textless{} b) form a fundamental system of neighbourhoods of b. Show that \(X \times X\) is completely normal but that X is not perfectly normal.
\item
  Show that if \(\mathbf{X}\) is a compact space such that \(\mathbf{X} \times \mathbf{X} \times \mathbf{X}\) is completely normal, then \(\mathbf{X}\) is metrizable (use Theorem 1 of §2, no. 4).
\end{enumerate}

¶ 13) Let \((\mathbf{X}_{t})_{t\in I}\) be an uncountable family of Hausdorff topological spaces, each containing more than one point. For each \(t\in I\) , let \(a_{t}\) and \(b_{t}\) be two distinct points of \(\mathbf{X}_{t}\) . Let \(\mathbf{X}\) be the product space \(\prod_{t\in I}\mathbf{X}_{t}\) , and let Y be the subspace of X consisting of these points \((x_{t})_{t\in I}\) such that \(x_{t}=a_{t}\) for all but a countable set of indices t; and let b be the point \(b_{t}\) of X.

In the product space \(\mathbf{X} \times \mathbf{Y}\) , let \(\mathbf{A}\) be the diagonal of \(\mathbf{Y} \times \mathbf{Y}\) and let \(\mathbf{B}\) be the set \(\{b\} \times \mathbf{Y}\) . Show that \(\mathbf{A}\) and \(\mathbf{B}\) are closed sets and that there is no pair of open sets \(\mathbf{U}, \mathbf{V}\) in \(\mathbf{X} \times \mathbf{Y}\) such that \(\mathbf{A} \subset \mathbf{U}, \mathbf{B} \subset \mathbf{V}\) , and \(\mathbf{U} \cap \mathbf{V} = \emptyset\) . {[}Let \(\mathbf{U}, \mathbf{V}\) be two open sets such that \(\mathbf{A} \subset \mathbf{U}\) and \(\mathbf{B} \subset \mathbf{V}\) ; show that there is an increasing sequence \((\mathbf{H}_n)\) of countable subsets of \(\mathbf{I}\) , such that if \(x_n\) denotes the point of \(\mathbf{Y}\) for which each coordinate with index \(i \in \mathbf{H}_n\) is equal to \(b_i\) , and each coordinate with index \(i \notin \mathbf{H}_n\) is equal to \(a_i\) , then the point \((x_{n+1}, x_n)\) belongs to \(\mathbf{V}\) for each integer \(n\) ; prove that the sequence \((x_{n+1}, x_n)\) tends to a point of \(\mathbf{A}\) , and deduce that \(\mathbf{U} \cap \mathbf{V} \neq \emptyset\) .{]}

Hence construct an example of a connected topological ring which is not normal.

Deduce also that the product of an uncountable family of Hausdorff topological spaces, each of which has at least two distinct points, cannot be completely normal (use the result above when each \(X_{t}\) is a discrete space of two points).

\begin{enumerate}
\def\labelenumi{\arabic{enumi})}
\setcounter{enumi}{13}
\item
  Let X be a non-normal Hausdorff space, and let A, B be two disjoint closed subsets of X such that there is no pair of disjoint open sets U, V in X satisfying the relations \(A \subset U\) and \(B \subset V\) . Let R be the equivalence relation on X whose equivalence classes are the set A, the set B, and the sets \(\{x\}\) where \(x \in \complement(A \cup B)\) . Show that the graph of R is closed in \(X \times X\) and that the relation R is closed, but that the quotient space X/R is not Hausdorff (cf.~Chapter I, § 8, no. 3, Proposition 8).
\item
  \begin{enumerate}
  \def\labelenumii{\alph{enumii})}
  \tightlist
  \item
    Let X be a normal space and let R be a closed equivalence relation on X. Show that the quotient space X/R is normal (use Proposition 10 of Chapter I, § 5, no. 4).
  \end{enumerate}
\end{enumerate}

\begin{enumerate}
\def\labelenumi{\alph{enumi})}
\setcounter{enumi}{1}
\item
  Let X be a locally compact σ-compact space, and let R be an equivalence relation on X whose graph is closed in X×X. Show that X/R is normal (cf.~Chapter I, § 10, Exercise 19).
\item
  Let X be the complement in R of the set of points of the form 1/n, where n is any integer except 0 or ±1. Consider the equivalence relation R on the metrizable space X for which the equivalence class of each \(x \in X\) , which is not an integer, consists of the points x and 1/x, and the equivalence class of each integer consists of that integer alone. Show that R is open and that the graph of R is closed in \(X \times X\) , but that X/R is not regular.
\end{enumerate}

\begin{enumerate}
\def\labelenumi{\arabic{enumi})}
\setcounter{enumi}{15}
\tightlist
\item
  For each covering \(\Re\) of a topological space \(X\) , let \(V_{\Re}\) denote the union of all the sets \(U \times U\) such that \(U \in \Re\) . A covering \(\Re\) of \(X\) is said to be even if there is a neighbourhood \(W\) of the diagonal \(\Delta\) of \(X \times X\) such that the covering \((W(x))_{x \in X}\) is finer than \(\Re\) . \(\Re\) is said to be divisible if there is a neighbourhood \(W\) of \(\Delta\) in \(X \times X\) such that \(\dot{W} \subset V_{\Re}\) . a) Show that every even open covering of \(X\) is divisible {[}cf.~Exercise 19 b){]}.
\end{enumerate}

\begin{enumerate}
\def\labelenumi{\alph{enumi})}
\setcounter{enumi}{1}
\tightlist
\item
  Let \(\Re\) be an open covering of \(X\) such that there exists a locally finite closed covering \(\mathfrak{S}\) of \(X\) which is finer than \(\Re\) . Show that \(\Re\) is even. (For every set \(A \in \mathfrak{S}\) , let \(U_A\) be a set of \(\Re\) which contains \(A\) , and let \(V_A\) be the union in \(X \times X\) of \(U_A \times U_A\) and \((X - A) \times (X - A)\) ; consider the set \(W = \bigcap_{A \in \mathfrak{S}} V_A\) .)
\end{enumerate}

¶ 17) a) If every finite open covering of a topological space X is divisible (Exercise 16), then X satisfies Axiom \((\mathrm{O}_{\mathbf{V}}^{\prime})\) (consider a binary open covering of X, i.e., one formed by two subsets of X). Conversely, if X satisfies \((\mathrm{O}_{\mathbf{V}}^{\prime})\) , then every finite open covering of X is even {[}use Theorem 3 of no. 3, which is valid provided that X satisfies \((\mathrm{O}_{\mathbf{V}}^{\prime})\) , to reduce to the case of a binary covering, and then again to deal with this case{]}. As R runs through the set of finite open coverings of X, the sets \(V_{R}\) form a fundamental system of entourages of a uniformity on X, compatible with the topology of X; this uniformity is called the ``uniformity of finite open coverings''.

\begin{enumerate}
\def\labelenumi{\alph{enumi})}
\setcounter{enumi}{1}
\item
  Suppose that X is normal. Show that the uniformity U of finite open coverings is the same as the coarsest uniformity \(U'\) for which all continuous mappings of X into {[}0, 1{]} are uniformly continuous {[}i.e., the uniformity induced on X by that of its Stone-Čech compactification (§ 1, Exercise 7){]}. {[}To show that U is coarser than \(U'\) , use Theorem 3 of no. 3 and Axiom \((\mathrm{O_V})\) in order to construct a finite family of pseudometrics on \(\mathbf{X}\) of the form \((x, y) \to |f(x) - f(y)|\) , in such a way that an entourage of the uniformity defined by these pseudometrics is contained in an entourage of the uniformity \(\mathfrak{U}]\) .
\item
  Let X be a normal space and let \(\beta\mathbf{X}\) be its Stone-Čech compactification. For each closed subset Y of X, let \(\overline{\mathbf{Y}}\) be the closure of Y in \(\beta\mathbf{X}\) . Show that the continuous mapping of \(\beta\mathbf{Y}\) into \(\overline{\mathbf{Y}}\) which extends the identity mapping of Y ( \(\beta\mathbf{Y}\) being the Stone-Čech compactification of Y) is a homeomorphism {[}either use \(b\) ) or else Theorem 2 and the universal property of \(\beta\mathbf{Y}]\) . If Y and Z are two closed subsets of X, show that \(\overline{\mathbf{Y}} \cap \overline{\mathbf{Z}} = \overline{\mathbf{Y} \cap \mathbf{Z}}\) in \(\beta\mathbf{X}\) (if \(x_0 \notin \overline{\mathbf{Y}} \cap \overline{\mathbf{Z}}\) and \(x_0 \in \overline{\mathbf{Y}}\) , consider a closed neighbourhood V of \(x_0\) in \(\beta\mathbf{X}\) , such that V \(\cap\) Y and Z are disjoint closed subsets of X).
\end{enumerate}

¶ 18) A family \((A_{\alpha})\) of subsets of a topological space X is said to be discrete if, for each \(x \in X\) , there is a neighbourhood of x which meets at most one set \(A_{\alpha}\) of the family. A Hausdorff space X is said to be collectively normal if, for each discrete family \((A_{\alpha})\) of closed subsets of X, there is a family \((U_{\alpha})\) of mutually disjoint open subsets of X such that \(A_{\alpha} \subset U_{\alpha}\) for each index \(\alpha\) . Every collectively normal space is normal.

\begin{enumerate}
\def\labelenumi{\alph{enumi})}
\item
  Every open covering of a completely regular space X is divisible (Exercise 16) if and only if the set of neighbourhoods of the diagonal \(\Delta\) in \(\mathbf{X} \times \mathbf{X}\) is the filter of entourages of the universal uniformity of X ( \(\S\) 1, Exercise 5). Show that if this condition is satisfied, then X is collectively normal {[}if \((\mathbf{A}_{\alpha})\) is a discrete family of closed subsets of X, consider the open covering \((\mathbf{V}_{\alpha})\) , where \(\mathbf{V}_{\alpha} = \mathbf{X} - \bigcup_{\beta \neq \alpha} \mathbf{A}_{\beta}\) for each \(\alpha\) {]}.
\item
  Let Y = {[}a, b{]} be the locally compact space defined in Chapter I, § 9, Exercise 12; let Y₀ be the set Y endowed with the discrete topology, and let Z be the set Y₀ ∪ \{b\}, where the points of Y₀ are open sets, and the sets {]}x, b{[}, where x runs through Y₀, form a fundamental system of neighbourhoods of b. Show that the space X = Y × Z is collectively normal (recall that no infinite family of subsets of Y can be discrete, and use Exercise 4). Let R be the open covering of X consisting of Y × Y₀ and the products {[}a, x{]} × {]}x, b{]} (where x \textless{} b); show that R is not divisible and that consequently the set of neighbourhoods of the diagonal Δ in X × X is not the filter of entourages of a uniformity on X {[}use Exercise 12 a) of Chapter I, § 9{]}.
\end{enumerate}

¶ 19) a) A regular space X is paracompact if and only if every open covering of X is even (Exercise 16). {[}To prove necessity, use Lemma 5 and Exercise 16 b); for sufficiency, use Exercise 16 a) of § 4, Proposition 2 of § 1, no. 4, and Theorem 4 of § 4, no. 5{]}.

\begin{enumerate}
\def\labelenumi{\alph{enumi})}
\setcounter{enumi}{1}
\item
  Give an example of a divisible non-even covering of a collectively normal, non-paracompact space {[}use a) and Chapter II, § 4, Exercise 4{]}.
\item
  Show that a paracompact space X is complete with respect to its universal uniformity (§ 1, Exercise 5). {[}If a Cauchy filter F with respect to this uniformity has no cluster point, then every point \(x \in X\) has an open neighbourhood \(V_{x}\) which does not meet at least one set of F; use a) and Exercise 18 a).{]}
\end{enumerate}

\begin{enumerate}
\def\labelenumi{\arabic{enumi})}
\setcounter{enumi}{19}
\tightlist
\item
  \begin{enumerate}
  \def\labelenumii{\alph{enumii})}
  \tightlist
  \item
    A regular space X is paracompact if, for each open covering R of X, there is a sequence \((\mathfrak{S}_{n})\) such that \(\mathfrak{S} = \bigcup_{n} \mathfrak{S}_{n}\) is an open covering of \(\mathbf{X}\) which refines \(R\) and such that each \(\mathfrak{S}_n\) is a locally finite family (cf.~Lemmas 4, 5 and 6).
  \end{enumerate}
\end{enumerate}

\begin{enumerate}
\def\labelenumi{\alph{enumi})}
\setcounter{enumi}{1}
\item
  Deduce from a) that in a paracompact space every countable union of closed sets is a paracompact subspace (cf.~§ 5, Exercise 15).
\item
  In a regular space X, let F be a locally finite family of closed sets each of which is a paracompact subspace of X. Show that the union of the sets of F is a paracompact subspace of X (use Lemma 6 of no. 5; cf.~§ 5, Exercise 15).
\item
  Show that if X is a paracompact space and Y is a regular space which is a countable union of compact sets, then \(X \times Y\) is paracompact {[}embed Y in its Stone-Čech compactification and use b); cf.~§ 5, Exercise 16{]}.
\end{enumerate}

¶ 21) a) Let X be a paracompact space and let R be a closed equivalence relation on X such that every equivalence class with respect to R is compact. Show that X/R is paracompact {[}use Exercise 15 a), Theorem 3 of no. 3, and Lemma 6 of no. 5; cf.~Exercise 15 c){]}.

\begin{enumerate}
\def\labelenumi{\alph{enumi})}
\setcounter{enumi}{1}
\item
  Let X be a Hausdorff space and let R be a closed equivalence relation on X such that every equivalence class with respect to R is compact. Show that if X/R is paracompact, then so is X (use the argument of Chapter I, § 9, no. 10, Proposition 17).
\item
  Let Y be a locally compact but not paracompact space (cf.~Chapter I, § 9, Exercise 12), and let \(Y_{0}\) be the one-point compactification of Y. Let X be the topological sum of Y and \(Y_{0}\) , and let R be the equivalence relation on X which identifies each point of Y with its canonical image in \(Y_{0}\) . The relation R is open, every equivalence class mod R contains at most two points, and X/R is homeomorphic to \(Y_{0}\) , hence is compact.
\end{enumerate}

\begin{enumerate}
\def\labelenumi{\arabic{enumi})}
\setcounter{enumi}{21}
\item
  A regular space \(\mathbf{X}\) is metrizable if and only if there exists a sequence \((\mathfrak{B}_n)\) of locally finite families of open subsets of \(\mathbf{X}\) such that \(\mathfrak{B} = \bigcup_{n} \mathfrak{B}_n\) is a base of the topology of \(\mathbf{X}\) . {[}To show that the condition is necessary, use Lemma 3 of no. 5. To show that it is sufficient, show first that X is paracompact by using Lemmas 4, 5 and 6. For each pair of integers m, n and each \(U \in \mathfrak{B}_{m}\) , let \(U'\) be the union of the sets \(V \in \mathfrak{B}_{n}\) such that \(\overline{V} \subset U\) . Show that \(\overline{U}' \subset U\) . Let \(f_{U}: X \to [0, 1]\) be a continuous mapping which is equal to 0 on X---U and to 1 on \(U'\) , and let \(d_{mn}(x, y) = \sum_{U \in \mathfrak{B}_{m}} |f_{U}(x) - f_{U}(y)|\) ; show that the pseudometrics \(d_{mn}\) define the topology of X.{]} (``Theorem of Nagata-Smirnov.''). In particular, every regular space with a countable base is metrizable.
\item
  \begin{enumerate}
  \def\labelenumii{\alph{enumii})}
  \tightlist
  \item
    Every regular Lindelöf space (Chapter I, § 9, Exercise 15) is paracompact {[}use Exercise 20 a){]}.
  \end{enumerate}
\end{enumerate}

\begin{enumerate}
\def\labelenumi{\alph{enumi})}
\setcounter{enumi}{1}
\tightlist
\item
  The product of a Lindelöf space and a compact space is a Lindelöf space (argue as in the proof of Proposition 17 of Chapter I, § 9, no. 10; cf.~§ 5, Exercise 16).
\end{enumerate}

\begin{enumerate}
\def\labelenumi{\arabic{enumi})}
\setcounter{enumi}{23}
\tightlist
\item
  \begin{enumerate}
  \def\labelenumii{\alph{enumii})}
  \tightlist
  \item
    Let X be a metrizable space and let R be a closed equivalence relation on X, all of whose equivalence classes are compact. Show that X/R is metrizable. {[}Apply the Nagata-Smirnov theorem (Exercise 22) by proving the analogue of Lemma 3 of no. 5 for open coverings of X formed of sets which are saturated with respect to R; let \(F_{n\alpha}\) denote the largest saturated open set contained in the set of all \(x \in X\) such that \(d(x, X - U_{\alpha}) > 2^{-n}\) , let \(F_{n\alpha}'\) denote the saturation of the set of all \(x \in X\) such that \(d(x, X - U_{\alpha}) \geqslant 2^{-n}\) and let \(G_{n\alpha}\) denote the set of all \(x \in F_{n\alpha}\) such that \(x \notin F_{n+1,\beta}'\) for all \(\beta < \alpha\) .{]}
  \end{enumerate}
\end{enumerate}

\begin{enumerate}
\def\labelenumi{\alph{enumi})}
\setcounter{enumi}{1}
\item
  Extend the result of a) to the case in which the relation R is closed and every point of X/R has a countable fundamental system of neighbourhoods {[}use Exercise 20 b) of § 2{]}.
\item
  Let X be a topological space and let \((\mathbf{F}_{\alpha})\) be a locally finite closed covering of X. Show that if each of the subspaces \(F_{\alpha}\) is metrizable then X is metrizable (consider X as a quotient space of the topological sum of the \(F_{\alpha}\) ).
\item
  A paracompact space, in which every point has a metrizable neighbourhood, is metrizable {[}apply c); cf.~Chapter I, § 9, Exercise 12 b), also Chapter IX, § 4, Exercise 25 d), § 2, Exercise 15 and § 5, Exercise 15{]}.
\end{enumerate}

\begin{enumerate}
\def\labelenumi{\arabic{enumi})}
\setcounter{enumi}{24}
\tightlist
\item
  A Hausdorff space X is said to be metacompact if, given any open covering R of X, there is a point-finite open covering of X which refines R.
\end{enumerate}

\begin{enumerate}
\def\labelenumi{\alph{enumi})}
\item
  Show that every closed subspace of a metacompact space is metacompact, and that if every open subspace of a metacompact space X is metacompact, then every subspace of X is metacompact.
\item
  Let X be a Hausdorff space and let R be a closed equivalence relation on X, all of whose equivalence classes are compact. Show that if X/R is metacompact, then X is metacompact {[}argue as in Exercise 21 b){]}.
\item
  Show that a topological space which is both metacompact and countably compact (§ 2, Exercise 14) is compact. {[}Use Zorn's lemma to show that every point-finite open covering contains a minimal open covering, and use § 2, Exercise 14 a).{]}
\end{enumerate}

\(d)\) The locally compact, completely normal space \(\mathbf{X} = [a, b[\) defined in Chapter I, § 9, Exercise 12 is not metacompact (*).

\begin{enumerate}
\def\labelenumi{\alph{enumi})}
\setcounter{enumi}{4}
\tightlist
\item
  Let Y be a metrizable space and let A be a subset of Y such that both A and its complement are dense in Y. Let X denote the non-regular space obtained by endowing Y with the topology generated by the open sets of Y and the set A. Show that X is metacompact. (For each \(x \in X\) , consider a neighbourhood \(V_{x}\) of x contained in a set of a given open covering R of X, such that \(V_{x} \subset A\) if \(x \in A\) , and such that \(V_{x}\) is a neighbourhood of x in Y if \(x \notin A\) . Note that the subspace A of X is paracompact, and that the union of the \(V_{x}\) for \(x \notin A\) is also paracompact.)
\end{enumerate}

\begin{enumerate}
\def\labelenumi{\arabic{enumi})}
\setcounter{enumi}{25}
\tightlist
\item
  \begin{enumerate}
  \def\labelenumii{\alph{enumii})}
  \tightlist
  \item
    Show that every pseudo-compact normal space X (§ 1, Exercise 22) is countably compact (§ 2, Exercise 14). {[}Show that if \((x_n)\) is a sequence of distinct points of X with no cluster point, then there is a continuous finite real-valued function f on X such that \(f(x_n) = n\) for all n.{]}
  \end{enumerate}
\end{enumerate}

\begin{enumerate}
\def\labelenumi{\alph{enumi})}
\setcounter{enumi}{1}
\tightlist
\item
  Let \(Y = [a, b[\) be the locally compact space defined in Chapter I, § 9, Exercise 12, and let \(Y_{0}\) be the compact space obtained by adjoining to Y a point at infinity (which may be identified with b). Let c be the smallest element of Y such that {[}a, c{[} is infinite, and let Z be the compact subspace {[}a, c{]} of Y. Let X be the complement of the point (b, c) in the compact product space \(Y_{0} \times Z\) . Show that X is locally compact but not normal {[}cf.~Exercise 12 a){]}; that X is pseudo-compact but not countably compact. Give an example of a closed subspace of X which is not pseudo-compact. Also, give an example of a lower semi-continuous function on X which does not attain its greatest lower bound {[}cf.~§ 2, Exercise 14 i){]}.
\end{enumerate}

\begin{enumerate}
\def\labelenumi{\arabic{enumi})}
\setcounter{enumi}{26}
\tightlist
\item
  A topological space X is said to be locally paracompact if every point of X has a closed neighbourhood which is a paracompact subspace of X.
\end{enumerate}

\begin{enumerate}
\def\labelenumi{\alph{enumi})}
\item
  Show that a locally paracompact space is completely regular.
\item
  In the compact space \(Y_{0} \times Z\) defined in Exercise 26 b), let H be the complement of the set of points \((b, y)\) where y \textless{} c.~Show that H is normal but not locally paracompact {[}cf.~Chapter I, § 9, Exercise 12 b){]}.
\item
  Let X be a topological space and let \(\Phi\) be a covering of X by closed sets which are paracompact subspaces of X, such that each \(A \in \Phi\) has a neighbourhood in X which belongs to \(\Phi\) . Let \(X'\) be the set which is the sum of X and a point \(\omega\) ; define a topology on \(X'\) by taking as a fundamental system of neighbourhoods of \(x \in X\) in \(X'\) the set of all neighbourhoods of x in X, and as a fundamental system of neighbourhoods of \(\omega\) the filter base generated by the complements in \(X'\) of the sets of \(\Phi\) {[}cf.~Exercise 20 c){]}. Show that \(X'\) is paracompact.
\end{enumerate}

\begin{enumerate}
\def\labelenumi{\arabic{enumi})}
\setcounter{enumi}{27}
\tightlist
\item
  Let X be a metric space, let d be the metric on X, let A be a closed subset of X and let \(f: A \to [1, 2]\) be a continuous real-valued function. Show that the real-valued function g on X which is such that \(g(x) = f(x)\) if \(x \in A\) , and
\end{enumerate}

\[
g (x) = \frac {\mathrm{1}}{d (x , \mathrm{A})} \inf _ {y \in \mathrm{A}} (f (y) d (x, y))
\]

if \(x \in \mathbb{C}A\) , is continuous on \(X\) . Hence give a proof of Theorem 2 for metric spaces.

\begin{enumerate}
\def\labelenumi{\arabic{enumi})}
\setcounter{enumi}{28}
\item
  Let X be a Hausdorff topological space. Then X is normal if and only if, for each closed subset A of X and each real-valued function f defined on A and continuous with respect to the induced topology, there exists a continuous pseudometric d on X with respect to which f is uniformly continuous. (To show that the condition is sufficient, show that every real-valued function f defined on a closed subset A, and continuous with respect to the induced topology, can be extended to a continuous function on X; for this, consider the Hausdorff space \(\overline{X}\) associated with the uniformity defined on X by the pseudometric d.)
\item
  Let X be a normal space and let g (resp. f) be an upper (resp. lower) semi-continuous function on X, such that \(g \leqslant f\) . Show that there is a continuous real-valued function h on X such that \(g \leqslant h \leqslant f\) . {[}Reduce to the case where f and g take their values in {[}0, 1{]}. For each dyadic number \(r \in [0, 1]\) , let \(\mathrm{F}(r)\) {[}resp. \(\mathrm{G}(r)\) {]} be the set of all \(x \in X\) such that \(f(x) \leqslant r\) (resp. \(g(x) < r\) ). Define by induction a family of open sets \(\mathrm{U}(r)\) (where r runs through the set of dyadic numbers contained in {[}0, 1{]}) such that when \(r < r'\) we have \(\mathrm{F}(r) \subset \mathrm{U}(r')\) , \(\overline{\mathrm{U}(r)} \subset \mathrm{U}(r')\) and \(\overline{\mathrm{U}(r)} \subset \mathrm{G}(r')\) , by imitating the proof of Theorem 1 of no. 1; then complete the proof as in Theorem 1.{]}
\end{enumerate}
% semantic-fragments: legacy-input-seam
\relax{}
